\section{研究目的}

% ＜改めて背景を洗う + 若い世代である必要性＞
ここまでに述べたように，自己開示は人間関係の親密化において重要な役割を果たす．特に，段階的な自己開示のプロセスとそれに対する適切な受容は，関係性の親密化に不可欠な要素である．しかし，現代では否定や拒絶への不安から深い自己開示を避ける傾向にあり，人間関係の希薄化につながっている．
% ＜音楽という新しいアプローチ＞
この課題に対し，1.4節で述べたように，音楽嗜好の共有は有効なアプローチとなる可能性がある．音楽は個人の価値観や感情を反映する要素であり，その共有をきっかけとして音楽にまつわる経験や個人的な価値観を語り合うことで，自然な形での自己開示が期待できる．

% ＜目的＞
以上より，本研究では，音楽嗜好の共有を通じた自己開示を促進することで，関係の親密化を目指す．